% !PRINTED 143
\weberpage{143}

One can also continue the system of formulae (8) in the opposite direction, if one forms the equations
\[
 S[\alpha^{-1}f(\alpha)]=0,\qquad
 S[\alpha^{-2}f(\alpha)]=0,\ldots .
\]
In this way one obtains a means of expressing the power sums $s_\nu$ also for negative exponents $\nu$ by the fundamental functions. These sums of negative powers likewise belong to the symmetric functions, although no longer to the \gesperrt{whole}, but to the \gesperrt{fractional}, ones. They pass into whole functions only after multiplication by powers of the product $\alpha_1\alpha_2\cdots\alpha_n$. For completeness we put down here the first two of these formulae:
\begin{equation}
\begin{aligned}
0&=s_{n-1}+a_1s_{n-2}+a_2s_{n-3}+\cdots+n a_{n-1}+a_ns_{-1},\\
0&=s_{n-2}+a_1s_{n-3}+a_2s_{n-4}+\cdots+a_{n-1}s_{-1}+a_ns_{-2}.
\end{aligned}
\tag{11}
\end{equation}
which, by (7), can also be written as
\[
 a_{n-1}+a_ns_{-1}=0,\qquad 2a_{n-2}+a_{n-1}s_{-1}+a_ns_{-2}=0.
\]

\begin{center}
{\large \S.\ 43.}

\vspace{0.5em}
{\large\gesperrt{Proof of the Principal Theorem for Two Variables.}}
\end{center}

We now pass to the proof of the fundamental theorem of the theory of symmetric functions, that they are all rationally expressible by the symmetric fundamental functions. Since we use complete induction as the means of proof, we first derive the theorem under the assumption that only two independent variables $\alpha,\beta$ are given, but in a way which at the same time will bring out the guiding thought for the general proof.

We denote the symmetric fundamental functions by
\begin{equation}
 a=-(\alpha+\beta),\qquad b=\alpha\beta
\tag{1}
\end{equation}
and accordingly put
\begin{equation}
 f(x)=(x-\alpha)(x-\beta)=x^2+ax+b.
\tag{2}
\end{equation}

Let $S(\alpha,\beta)$ be any whole rational and symmetric function of $\alpha$ and $\beta$. We can substitute for $\beta$ from (1) the value $-(a+\alpha)$, and obtain, if we arrange according to powers of $\alpha$,
\begin{equation}
 S(\alpha,\beta)=S(\alpha,-a-\alpha)=A_0\alpha^m+A_1\alpha^{m-1}+\cdots+A_{m-1}\alpha+A_m,
\tag{3}
\end{equation}
%CONT>>
